\documentclass[11pt]{article}
\usepackage[margin=1in]{geometry}
\usepackage{amsmath,amssymb,amsthm,mathtools}
\usepackage{graphicx}
\usepackage{microtype}
\usepackage{enumitem}
\usepackage{xcolor}
\usepackage[hidelinks]{hyperref}
\setlength{\emergencystretch}{3em}

\newtheorem{theorem}{Theorem}
\newtheorem{corollary}[theorem]{Corollary}
\newtheorem{lemma}[theorem]{Lemma}
\newtheorem{proposition}[theorem]{Proposition}
\theoremstyle{definition}
\newtheorem{example}[theorem]{Example}
\newtheorem{remark}[theorem]{Remark}

\newcommand{\R}{\mathbb R}
\newcommand{\cC}{\mathcal C}
\newcommand{\cE}{\mathcal E}
\newcommand{\cH}{\mathcal H}
\newcommand{\cI}{\mathcal I}
\newcommand{\phib}{\phi_{\mathrm{bi}}}
\newcommand{\pJ}{p_J}
\newcommand{\eps}{\varepsilon}
\newcommand{\domm}{\operatorname{dom}}
\newcommand{\conv}{\operatorname{conv}}

\title{\textbf{Canonical Linear Core, Generated Hull, and Duality}\\
for Non-Even Orlicz Cones}
\author{Candidate major partial result for arXiv:2204.12282}
\date{12 August 2026}

\begin{document}
\maketitle

\begin{center}
\fbox{\parbox{0.92\textwidth}{
\textbf{Status.} Candidate major partial, likely valid; human review requested.
The broad program of a complete nonseparable Orlicz-cone theory and a
systematic calculus directly on locally convex cones remains open.
}}
\end{center}

\section{Source question and scope}

The source is Thomas Ruf, \emph{Non-Separably Valued Orlicz Spaces I},
arXiv:2204.12282v2 (2023) \cite{Ruf2023}.  On PDF page~3 the paper asks:
\begin{quote}
\emph{Can an analogue theory be developed for Orlicz cones generated by a
non-even convex Orlicz integrand $\phi\colon\Omega\times X\to[0,\infty]$?}
\end{quote}
It then asks for an associated subdifferential calculus on locally convex
cones, motivated by unilateral constraints.  Figure~\ref{fig:question} is a
crop of the source page.

\begin{figure}[ht]
\centering
\includegraphics[width=\textwidth]{figures/open_problem_crop.png}
\caption{The open and follow-up questions in arXiv:2204.12282v2, p.~3.}
\label{fig:question}
\end{figure}

This packet gives a canonical symmetric core, a canonical generated linear
hull valid without separability, exact cone-dual extension and conjugate
formulas, an attained-subdifferential rule, and a pointwise Banach-Orlicz
realization under finite-dimensional normal-integrand hypotheses.  It does
not claim the full program.

\section{Setting}

Let $(\Omega,\mathcal A,\mu)$ be a measure space, let $X$ be a real Banach
space, and let $\cE$ be the vector space of a.e.\ classes of strongly
measurable maps $\Omega\to X$.  Assume that
\[
 \phi\colon\Omega\times X\longrightarrow[0,\infty]
\]
is integrally separably measurable in the sense of \cite{Ruf2023}, and that
for a.e.\ $\omega$ the function $\phi_\omega=\phi(\omega,\cdot)$ is proper,
lower semicontinuous and convex, satisfies $\phi_\omega(0)=0$, is continuous
at $0$, and has bounded sublevel sets.  We do \emph{not} assume evenness.
Define
\[
 I(u):=\int_\Omega\phi(\omega,u(\omega))\,d\mu,\qquad
 p(u):=\inf\{\alpha>0:I(u/\alpha)\le1\},
\]
with the usual extended-value conventions, and put
\[
 \cC:=\{u\in\cE:p(u)<\infty\}.
\]
Convexity at the origin makes $p$ an extended sublinear gauge, so $\cC$ is a
convex cone.  The elementary modular inequalities used below are
\[
 p(z)\le1\Longrightarrow I(z)\le p(z),\qquad
 I(z)\le M<\infty\Longrightarrow p(z)\le\max\{1,M\}.
\tag{2.1}\label{eq:modular-gauge}
\]
The first follows by radial left-continuity; the second follows by convex
rescaling.

\section{Idea of the proof}

There are two different symmetric objects hidden in a one-sided Orlicz cone.
To stay \emph{inside} the cone while allowing signs, a vector and its negative
must both be admissible.  This forces the reflected maximum
\[
 \phib(\omega,x)=\max\{\phi(\omega,x),\phi(\omega,-x)\}.
\]
It yields the largest linear subspace contained in the cone.

To generate a linear space from the \emph{whole} cone, one instead writes
$u=v-w$ with $v,w\in\cC$.  Pointwise minimization over such decompositions
can fail to commute with integration in a nonseparable range.  The key
selection-free replacement is therefore the function-space infimal
convolution
\[
 J(u)=\inf_{v\in\cE}\{I(v)+I(v-u)\}.
\]
Its Luxemburg domain is exactly $\cC-\cC$.  Its symmetric gauge is equivalent
to the Grothendieck quotient norm obtained by minimizing $p(v)+p(w)$ over
$u=v-w$.  This quotient description makes the dual extension theorem
automatic, while expanding the Fenchel conjugate separates the forward and
backward terms exactly.  When finite-dimensional measurable selection is
available, $J$ becomes the integral of the pointwise infimal convolution and
the hull becomes an ordinary Banach Orlicz space.

\section{The maximal linear core}

\begin{theorem}[canonical maximal core]\label{thm:core}
Set
\[
 \phib(\omega,x):=\max\{\phi(\omega,x),\phi(\omega,-x)\}.
\]
Then $\phib$ is an even Orlicz integrand and
\[
 \cC\cap(-\cC)=L_{\phib}(\mu).
\tag{4.1}\label{eq:core}
\]
This is the unique maximal vector subspace contained in $\cC$.  On this
space, if $p_{\mathrm{bi}}$ is the Luxemburg norm generated by $\phib$, then
\[
 \max\{p(u),p(-u)\}
 \ \le\ p_{\mathrm{bi}}(u)\
 \le\ 2\max\{p(u),p(-u)\}.
\tag{4.2}\label{eq:corebounds}
\]
The universal constant $2$ cannot be decreased.
\end{theorem}

\begin{proof}
The maximum of two proper lower semicontinuous convex functions is again
lower semicontinuous and convex.  The definition makes $\phib$ even;
continuity at zero, coercivity, and integral separable measurability pass to
the maximum.  Thus the even theory of \cite{Ruf2023} applies.

Since $\phib\ge\phi$ and $\phib(\omega,x)\ge\phi(\omega,-x)$,
\[
 p(u)\le p_{\mathrm{bi}}(u),\qquad p(-u)\le p_{\mathrm{bi}}(u).
\]
This proves the left inequality in \eqref{eq:corebounds} and the inclusion
$L_{\phib}\subset\cC\cap(-\cC)$.

Conversely, let $u,-u\in\cC$ and take
$a>\max\{p(u),p(-u)\}$.  Then $I(u/a)\le1$ and $I(-u/a)\le1$.  Hence
\[
 I_{\mathrm{bi}}(u/a)
 \le I(u/a)+I(-u/a)\le2.
\]
Convexity at zero gives
$I_{\mathrm{bi}}(u/(2a))\le1$, so
$p_{\mathrm{bi}}(u)\le2a$.  Letting $a$ decrease proves the right inequality
and the reverse inclusion in \eqref{eq:core}.

The intersection $\cC\cap(-\cC)$ is a vector space by
\eqref{eq:core}.  If $V$ is any vector subspace contained in $\cC$, then
$u\in V$ implies $-u\in V\subset\cC$, so $V\subset\cC\cap(-\cC)$.
Optimality of $2$ is proved in Example~\ref{ex:sharp}.
\end{proof}

\begin{remark}
The norm $u\mapsto\max\{p(u),p(-u)\}$ is therefore complete on the core,
because it is equivalent to the Banach-Orlicz norm $p_{\mathrm{bi}}$.
\end{remark}

\section{The generated hull and its Grothendieck norm}

Define the even convex functional
\[
 J(u):=\inf_{v\in\cE}\{I(v)+I(v-u)\},
\qquad
 \pJ(u):=\inf\{\alpha>0:J(u/\alpha)\le1\}.
\tag{5.1}\label{eq:J}
\]
Also put
\[
 \cH:=\cC-\cC,\qquad
 r(u):=\inf_{\substack{v,w\in\cC\\u=v-w}}\{p(v)+p(w)\}.
\tag{5.2}\label{eq:r}
\]

\begin{theorem}[selection-free hull and sharp norm comparison]
\label{thm:hull}
The functional $J$ is even and convex, and
\[
 \{u\in\cE:\pJ(u)<\infty\}=\cH=\cC-\cC.
\tag{5.3}\label{eq:domainJ}
\]
Moreover, $p_J$ and $r$ are norms on $\cH$ and
\[
 \pJ(u)\le r(u)\le2\pJ(u)\qquad(u\in\cH).
\tag{5.4}\label{eq:hullbounds}
\]
The constant $2$ cannot be decreased, even for an even autonomous integrand
on a one-point measure space.
\end{theorem}

\begin{proof}
Writing $w=v-u$, formula \eqref{eq:J} is the infimal convolution of $I$
with its reflection.  It is convex, and
\[
 J(-u)=\inf_v\{I(v)+I(v+u)\}
      =\inf_w\{I(w-u)+I(w)\}=J(u).
\]

If $u=v-w$ with $v,w\in\cC$, choose a common scale $a>0$ at which
$I(v/a)$ and $I(w/a)$ are finite.  Then
$J(u/a)\le I(v/a)+I(w/a)<\infty$, and a further convex rescaling makes this
quantity at most $1$.  Thus $\pJ(u)<\infty$.
Conversely, if $\pJ(u)<\infty$, then $J(u/a)<\infty$ for some $a>0$.
There is a $v\in\cE$ for which both $I(v)$ and $I(v-u/a)$ are finite.
Putting $w=v-u/a$ gives $u=av-aw$ with $av,aw\in\cC$.  This proves
\eqref{eq:domainJ} without an attainment or measurable-selection assumption.

For the first inequality in \eqref{eq:hullbounds}, fix
$u=v-w$, take $a>p(v)$ and $b>p(w)$, and put $s=a+b$.  Convexity gives
\[
 I(v/s)\le\frac{a}{s}I(v/a)\le\frac{a}{s},\qquad
 I(w/s)\le\frac{b}{s}I(w/b)\le\frac{b}{s}.
\]
Therefore $J(u/s)\le1$ and $\pJ(u)\le a+b$.  Take the infimum.

For the reverse inequality, take $a>\pJ(u)$, so $J(u/a)\le1$ after
decreasing an admissible scale if necessary.  For every $\eps>0$ there are
$v,w\in\cE$ with $v-w=u/a$ and
\[
 I(v)+I(w)\le1+\eps.
\]
Each term is at most $1+\eps$, hence convex rescaling gives
$p(v),p(w)\le1+\eps$.  Consequently
$r(u)\le2a(1+\eps)$.  Let $\eps\downarrow0$ and $a\downarrow\pJ(u)$.

It remains only to note separation.  If $r(u)=0$, choose
$u=v_n-w_n$ with
$p(v_n)+p(w_n)\le2^{-2n}$.  By \eqref{eq:modular-gauge},
\[
 I(2^n v_n)+I(2^n w_n)\le2^{1-n}.
\]
The sum of these integrals is finite.  Hence both integrands tend to zero
a.e.  Coercivity then forces $v_n\to0$ and $w_n\to0$ a.e.; otherwise
$2^n\|v_n(\omega)\|$ or $2^n\|w_n(\omega)\|$ would diverge along a
subsequence while the corresponding integrand tends to zero.  Thus
$u=v_n-w_n=0$ a.e.  The equivalence \eqref{eq:hullbounds} separates $p_J$
as well.  Optimality is in Example~\ref{ex:sharp}.
\end{proof}

\section{Cone duality and subdifferential splitting}

Call a map $A\colon\cC\to\R$ a cone-linear functional if it is additive and
positively homogeneous.  Give it the norm
\[
 \|A\|_{\cC}:=\sup\{|A(v)|:v\in\cC,\ p(v)\le1\}.
\]

\begin{theorem}[Grothendieck extension and exact bidirectional duality]
\label{thm:dual}
\begin{enumerate}[label=\textnormal{(\alph*)}]
\item Every cone-linear $A$ with $\|A\|_{\cC}<\infty$ has the unique
linear extension
\[
 \ell(v-w)=A(v)-A(w)\qquad(v,w\in\cC)
\tag{6.1}\label{eq:extension}
\]
to $\cH$.  It is continuous for $r$, and
$\|\ell\|_{r,*}=\|A\|_{\cC}$.

\item For a linear functional $\ell$ on $\cH$, define
\[
 I^*(\ell):=\sup_{v\in\cC}\{\ell(v)-I(v)\},\qquad
 J^*(\ell):=\sup_{u\in\cH}\{\ell(u)-J(u)\}.
\]
Then, with extended nonnegative values,
\[
 J^*(\ell)=I^*(\ell)+I^*(-\ell).
\tag{6.2}\label{eq:conjugate}
\]

\item Suppose $u=v-w$, $v,w\in\cC$, attains
$J(u)=I(v)+I(w)$.  With subdifferentials taken in any fixed separating
dual pair containing $\ell$,
\[
 \ell\in\partial J(u)
 \quad\Longleftrightarrow\quad
 \ell\in\partial I(v)\ \hbox{ and }\ -\ell\in\partial I(w).
\tag{6.3}\label{eq:subdiff}
\]
\end{enumerate}
\end{theorem}

\begin{proof}
If $v-w=v'-w'$, then $v+w'=v'+w$.  Additivity gives
$A(v)-A(w)=A(v')-A(w')$, so \eqref{eq:extension} is well-defined and
plainly linear.  For every decomposition,
\[
 |\ell(v-w)|\le\|A\|_{\cC}\{p(v)+p(w)\}.
\]
Thus $\|\ell\|_{r,*}\le\|A\|_{\cC}$.  Conversely $r(v)\le p(v)$ for
$v\in\cC$, so restriction gives the reverse norm inequality and uniqueness.

For the conjugate, parameterize $u=v-w$:
\begin{align*}
 J^*(\ell)
 &=\sup_{\substack{u\in\cH,\ v\in\cE}}
   \{\ell(u)-I(v)-I(v-u)\}\\
 &=\sup_{v,w\in\cC}
   \{\ell(v)-I(v)-\ell(w)-I(w)\}\\
 &=I^*(\ell)+I^*(-\ell).
\end{align*}
No minimax interchange or attainment is used.

For the last assertion, if the two right-hand subgradient inclusions hold,
add their subgradient inequalities for any $u'=v'-w'$ and then take the
infimum over decompositions of $u'$.  This gives
$J(u')\ge J(u)+\ell(u'-u)$.
Conversely, Fenchel equality for $\ell\in\partial J(u)$, the attained
decomposition, and \eqref{eq:conjugate} says
\[
 [I(v)+I^*(\ell)-\ell(v)]
 +[I(w)+I^*(-\ell)+\ell(w)]=0.
\]
Both brackets are nonnegative Fenchel gaps, so both vanish.  This is exactly
the pair of subgradient inclusions in \eqref{eq:subdiff}.
\end{proof}

\begin{corollary}[polar formula]\label{cor:polar}
Let
\[
 \mathfrak a_I(\ell):=
 \inf_{\alpha>0}\alpha\{1+I^*(\ell/\alpha)\}.
\]
Then
\[
 \|\ell\|_{r,*}
 =\sup_{p(v)\le1}|\ell(v)|
 =\max\{\mathfrak a_I(\ell),\mathfrak a_I(-\ell)\}.
\tag{6.4}\label{eq:polar}
\]
\end{corollary}

\begin{proof}
The $r$-unit ball is the absolutely convex hull of the one-sided unit ball
$\{p\le1\}$, in the sense of Minkowski gauges.  Its support function is
therefore $\sup_{p(v)\le1}|\ell(v)|$.  The one-sided
Luxemburg--Amemiya polar identity
\[
 \sup_{p(v)\le1}\ell(v)
 =\inf_{\alpha>0}\alpha\{1+I^*(\ell/\alpha)\}
\]
is the usual gauge-conjugacy calculation; the proof in
\cite[\S3.1]{Ruf2023} uses only convexity, $I(0)=0$, and radial
left-continuity for this one-sided statement.  Apply it to $\ell$ and
$-\ell$.
\end{proof}

\section{A pointwise Banach-Orlicz hull}

\begin{theorem}[finite-dimensional pointwise realization]
\label{thm:pointwise}
Assume additionally that $(\Omega,\mathcal A,\mu)$ is complete and
$\sigma$-finite, $X$ is finite-dimensional, and $\phi$ is a normal integrand.
Define
\[
 \psi(\omega,x):=
 \min_{y\in X}\{\phi(\omega,y)+\phi(\omega,y-x)\}.
\tag{7.1}\label{eq:psi}
\]
Then $\psi$ is an even Orlicz integrand,
\[
 J(u)=I_\psi(u)\quad(u\in\cE),\qquad
 \cC-\cC=L_\psi(\mu),\qquad
 \pJ=\|\cdot\|_\psi.
\tag{7.2}\label{eq:pointwise}
\]
In particular, $(\cC-\cC,\pJ)$ is a Banach space and the ordinary even
Orlicz theory applies to the generated hull.
\end{theorem}

\begin{proof}
For fixed $(\omega,x)$, coercivity makes
$y\mapsto\phi(\omega,y)+\phi(\omega,y-x)$ inf-compact in finite
dimensions, so the minimum is attained.  Standard inf-projection and
measurable-selection results for normal integrands
\cite[Chapter 14]{RockafellarWets} show that $\psi$ is a normal integrand
and provide measurable $\eps$-minimizers.

The infimal convolution of a convex function with its reflection is convex
and even.  It vanishes continuously at zero because
\[
 0\le\psi(\omega,x)
 \le\min\{\phi(\omega,x),\phi(\omega,-x)\}\longrightarrow0.
\]
If $\|x_n\|\to\infty$ while $\psi(\omega,x_n)$ stayed bounded, minimizing
$y_n$ would make both $\phi(\omega,y_n)$ and
$\phi(\omega,y_n-x_n)$ bounded.  Bounded sublevels would make both
sequences bounded, contradicting $x_n=y_n-(y_n-x_n)$.  Thus $\psi$ is an
even Orlicz integrand.

Every measurable $v$ gives the pointwise upper candidate
$\psi(\omega,u(\omega))\le
\phi(\omega,v(\omega))+\phi(\omega,v(\omega)-u(\omega))$, hence
$I_\psi(u)\le J(u)$.  Measurable $\eps$-minimizers, with an integrable
positive error weight on the $\sigma$-finite space, give the reverse
inequality after integration and $\eps\downarrow0$.  Thus $J=I_\psi$.
Now \eqref{eq:domainJ} gives the set identity, and the gauge identity is
immediate.  Completeness follows from the even Orlicz-space theorem of
\cite{Ruf2023}.
\end{proof}

\section{Sharpness examples}

\begin{example}[both constants are optimal]\label{ex:sharp}
\textbf{Core.}  Let $\Omega$ have two atoms of mass $1/2$, let $X=\R$, and
for $0<\eps<1$ put
\[
 \phi_\eps(t)=t_++\eps t_-,\qquad u=(1,-1).
\]
The modular is positively homogeneous, so
\[
 \max\{p(u),p(-u)\}=\frac{1+\eps}{2},\qquad
 p_{\mathrm{bi}}(u)=1.
\]
The ratio is $2/(1+\eps)\to2$.

\smallskip
\textbf{Hull.}  On a one-point space let
\[
 \phi_M(t)=M(|t|-1)_+,\qquad M>0.
\]
Then
\[
 p(t)=\frac{M|t|}{M+1},\qquad
 J(t)=M(|t|-2)_+,\qquad
 \pJ(t)=\frac{M|t|}{2M+1}.
\]
Here $p$ is already symmetric, so $r=p$.  Hence
\[
 \frac{r(t)}{\pJ(t)}=\frac{2M+1}{M+1}\longrightarrow2.
\]
\end{example}

\section{Verification and novelty bounds}

The accompanying script \texttt{code/verify\_discrete\_models.py} checks
$500$ random finite-atom power models for \eqref{eq:corebounds}, $500$
random scalar power models for the coefficient form of
\eqref{eq:conjugate}, and the two sharpness families above.  These checks are
sanity tests, not proof.

The bounded novelty search covered the run's cheap indexes and exact/close
arXiv searches for non-even Orlicz integrands, Orlicz cones, asymmetric
Orlicz or modular spaces, symmetrization, and infimal convolution.  Ruf's
2024 thesis repeats the questions \cite{Ruf2024}.  A 2025 paper develops
max-symmetrization for abstract quasi-modular pseudometrics
\cite{Majozi2025}, but the checked text does not contain the maximal
Orlicz-cone core, the generated hull, the Grothendieck dual extension, or the
subdifferential split proved here.  Novelty confidence is therefore
\emph{moderate}, not definitive: terminology in asymmetric modular theory is
fragmented.

\section{Limitations and review recommendation}

\begin{itemize}[leftmargin=*]
\item The result is a major partial answer, not a complete construction of
the source's requested theory on locally convex cones.
\item In the arbitrary nonseparable case, $J$ is deliberately defined at the
function-space level.  A pointwise integrand formula would require an
inf-projection and measurable-selection theorem that is unavailable under the
stated general hypotheses.
\item The subdifferential equivalence \eqref{eq:subdiff} requires attainment
of the optimal decomposition.  Without attainment, only approximate
variants are immediate.
\item Human review should focus on the extended-gauge conventions in
\eqref{eq:domainJ}, the one-sided polar identity in
Corollary~\ref{cor:polar}, and the normal-integrand selection step in
Theorem~\ref{thm:pointwise}.
\end{itemize}

\sloppy
\begin{thebibliography}{9}

\bibitem{Ruf2023}
T.~Ruf,
\emph{Non-Separably Valued Orlicz Spaces I},
arXiv:2204.12282v2, 2023.

\bibitem{Ruf2024}
T.~D.~Ruf,
\emph{Orlicz-type Function Spaces and Generalized Gradient Flows with
Degenerate Dissipation Potentials in Non-Reflexive Banach Spaces: Theory and
Application},
doctoral dissertation, arXiv:2401.17440, 2024.

\bibitem{RockafellarWets}
R.~T.~Rockafellar and R.~J.-B.~Wets,
\emph{Variational Analysis},
Grundlehren der mathematischen Wissenschaften 317, Springer, 1998.

\bibitem{Majozi2025}
P.~R.~Majozi,
\emph{On Quasi-Modular Pseudometric Spaces and Asymmetric Uniformities},
arXiv:2510.25200, 2025 (source version checked in 2026).

\end{thebibliography}
\fussy

\end{document}
